\pdfinfoomitdate=1
\pdftrailerid{}
\pdfsuppressptexinfo=15
\documentclass[10pt]{article}
\usepackage[a4paper,margin=23mm]{geometry}
\usepackage{amsmath,amssymb,amsthm,booktabs,array,etoolbox}
\usepackage[hidelinks]{hyperref}
\usepackage[T1]{fontenc}

\newtheorem{theorem}{Theorem}
\newtheorem{lemma}{Lemma}
\newtheorem{proposition}{Proposition}
\newcommand{\F}{\mathbb F}
\newcommand{\Z}{\mathbb Z}
\newcommand{\CW}{\operatorname{CW}}
\newcommand{\CGW}{\operatorname{CGW}}
\newcommand{\supp}{\operatorname{supp}}
\newcommand{\wt}{\operatorname{wt}}
\apptocmd{\thebibliography}{%
  \setlength{\itemsep}{0.4\baselineskip}%
}{}{}

\title{Nonexistence of a Circulant $\CGW(35,12;6)$\\
and the Case $\CW(105,36)$}
\author{Semyon Osadchii\\
\small Independent researcher\\
\small \texttt{Sposadchii@gmail.com}; \texttt{sp.osadchii@student.han.nl}}
\date{29 July 2026}

\begin{document}
\maketitle

\begin{abstract}
We give an exact computer-assisted proof that no circulant complex generalized
weighing matrix $\CGW(35,12;6)$ exists.  Reduction modulo the Eisenstein prime
$1-\omega$ places the support of a hypothetical first row in one of two
explicit ternary cyclic $[35,12]$ codes.  Complete enumeration gives 420
possible supports, and an exact affine-orbit calculation shows that they form
one orbit.  On a canonical support, reduction modulo $2$ yields a short
contradiction from the four correlations at shifts $2,4,8,10$.  The known
three-fold contraction of a hypothetical $\CW(105,36)$ would produce exactly
such a sixth-root sequence, so its nonexistence follows.  Python, standalone
C++, result-blind SageMath and result-blind GAP implementations reproduce the
finite classification by distinct routes; these computations are supporting
validation, while the final characteristic-two obstruction is algebraic.
\end{abstract}

\section{Statement and group-ring formulation}

Let $\mu_6=\langle\zeta_6\rangle$, where
$\zeta_6=e^{\pi i/3}=1+\omega$ and $\omega^2+\omega+1=0$.  A circulant
$\CGW(n,w;6)$ is a circulant matrix with entries in
$\{0\}\cup\mu_6$ whose product with its conjugate transpose is $wI_n$.
Equivalently, its first row
\[
 u=(u_0,\ldots,u_{n-1})\in(\{0\}\cup\mu_6)^n
\]
has Hamming weight $w$ and periodic Hermitian correlations
\[
 C_u(t)=\sum_{q\in\Z_n}u_q\overline{u_{q+t}}
 =\begin{cases}w,&t=0,\\0,&t\ne0.\end{cases}
\]

Put
\[
 U(X)=\sum_{q=0}^{34}u_qX^q
 \quad\text{in}\quad
 \Z[\omega][X]/(X^{35}-1),
\]
and let $\#$ conjugate coefficients and send $X$ to $X^{-1}$.  Thus a
circulant $\CGW(35,12;6)$ is exactly a weight-12 polynomial satisfying
\begin{equation}\label{eq:group-ring}
 UU^\#=12.
\end{equation}
No coordinate is prescribed to vanish.

\section{The factor-orbit reduction}

\begin{lemma}[Factor-orbit lemma]\label{lem:factor-orbit}
Let $F$ be a field, let $M\in F[X]$ be square-free, and suppose $\#$ induces
an involution on the irreducible factors of $M$.  If
\[
 vv^\#=0\pmod M,
\]
then every factor fixed by $\#$ divides $v$, and from every two-element orbit
$\{f,f^\#\}$ at least one of $f,f^\#$ divides $v$.
\end{lemma}

\begin{proof}
Each irreducible factor $f$ of the square-free polynomial $M$ is prime.
Therefore $f\mid vv^\#$ implies $f\mid v$ or $f\mid v^\#$.  The second
alternative is equivalent to $f^\#\mid v$.
\end{proof}

Reduce \eqref{eq:group-ring} modulo $\pi=1-\omega$.  Since
\[
 \Z[\omega]/(\pi)\cong\F_3
\]
and every element of $\mu_6$ maps to $1$ or $-1$, the reduction $V$ has
exactly the same support as $U$.  Moreover $12=0$ in $\F_3$, so
\begin{equation}\label{eq:f3-zero-product}
 V(X)V(X^{-1})=0\pmod{X^{35}-1}.
\end{equation}

The complete irreducible factorization over $\F_3$ is
\[
 X^{35}-1=(X-1)\Phi_5(X)\Phi_7(X)A(X)B(X),
\]
where $A$ and $B$ are reciprocal irreducible polynomials of degree 12.  In
low-to-high coefficient notation,
\[
\begin{aligned}
A&=[1,2,2,1,2,1,0,1,2,0,1,0,1],\\
B&=[1,0,1,0,2,1,0,1,2,1,2,2,1].
\end{aligned}
\]
All five displayed factors are distinct and irreducible, and their product is
$X^{35}-1$; this is therefore the complete square-free factorization.  The
first three factors are fixed by reciprocity, while $\{A,B\}$ is the sole
two-element reciprocal orbit.  Put $H=(X-1)\Phi_5\Phi_7$.  Lemma
\ref{lem:factor-orbit} gives $H\mid V$ and either $A\mid V$ or $B\mid V$.
Consequently the complete solution locus of
\eqref{eq:f3-zero-product} is the union of the two cyclic $[35,12]$ codes
generated by
\[
 G_A=(X-1)\Phi_5\Phi_7A,\qquad
 G_B=(X-1)\Phi_5\Phi_7B.
\]
Conversely, a multiple of $G_A$ has reciprocal divisible by $G_B$ (and
vice versa), so every word in either code satisfies
\eqref{eq:f3-zero-product}.
Their intersection is $\{0\}$: divisibility by both generators gives
$X^{35}-1\mid V$.  Since the reduction under consideration has support 12,
it is nonzero and lies in exactly one of the two codes.
Their coefficient vectors, again low-to-high, are
\[
\begin{aligned}
G_A={}&[2,0,1,0,1,1,0,2,0,1,2,1,1,1,0,2,0,1,1,1,2,2,1,1],\\
G_B={}&[2,2,1,1,2,2,2,0,1,0,2,2,2,1,2,0,1,0,2,2,0,2,0,1].
\end{aligned}
\]
The standalone publication audit verifies the factor product,
irreducibility, reciprocity, and both generator products exactly.

\section{Complete ternary enumeration and its certificate}

For each generator $G\in\{G_A,G_B\}$, form the 12 rows
$G,XG,\ldots,X^{11}G$ in $\F_3^{35}$.  Their leading coordinates are
distinct, so they have rank 12.  The enumeration visits every message in
$\F_3^{12}$, multiplies it by this generator matrix, computes the Hamming
weight, and retains the words of weight 12.  Direct periodic correlation is
checked on every retained word.  The exact counts are
\begin{center}
\begin{tabular}{lrr}
\toprule
&$\langle G_A\rangle$&$\langle G_B\rangle$\\
\midrule
messages enumerated&$3^{12}$&$3^{12}$\\
weight-12 codewords&420&420\\
distinct supports&210&210\\
\bottomrule
\end{tabular}
\end{center}
Each support occurs twice, for the scalar pair $v,-v$, and the two support
families are disjoint.

For an explicit coordinate anchor, the message
\[
 m=(1,0,0,2,0,1,0,0,0,0,0,0)
\]
in the $G_B$ branch gives the word
\[
\begin{split}
v={}&(2,2,1,2,0,0,0,2,0,0,1,0,2,0,0,0,2,0,0,0,0,\\
    &\hspace{20mm}1,1,0,0,0,1,0,1,0,0,0,0,0,0).
\end{split}
\]
It has zero periodic autocorrelation over $\F_3$ and support
\begin{equation}\label{eq:canonical-support}
 S=\{0,1,2,3,7,10,12,16,21,22,26,28\}.
\end{equation}

\begin{lemma}[Affine-orbit lemma]\label{lem:affine-orbit}
The 420 supports obtained above are exactly the orbit of $S$ under
\[
 q\longmapsto aq+b,\qquad
 a\in\Z_{35}^{\times},\quad b\in\Z_{35}.
\]
\end{lemma}

\begin{proof}
The exact enumeration serializes every support as a 35-bit integer, with bit
$q$ representing coefficient $q$.  Independently, all
$35\varphi(35)=840$ affine maps are applied to $S$.  The setwise stabilizer is
\[
 \{(a,b)=(1,0),(29,28)\},
\]
so the orbit contains $840/2=420$ supports.  Sorting both integer sets gives
byte-for-byte equality with no duplicates or omissions.  The complete ordered
list accompanies this paper.
\end{proof}

The numeric decimal-lines certificate has SHA-256
\[
\texttt{c067600256b37e077ee8a83889ec5e09347c397feea723466fb863ca2685f3d2}.
\]
For completeness, affine relabelling preserves the Hermitian correlation
condition directly.  If $a\in\Z_{35}^{\times}$, $b\in\Z_{35}$, and
$y_{aq+b}=\lambda x_q$ with $\lambda\overline{\lambda}=1$, then, writing
$C_h(x)=\sum_qx_q\overline{x_{q+h}}$, the change of variable
$q\mapsto aq+b$ gives
\[
 C_{ah}(y)=C_h(x).
\]
Multiplication by $a$ permutes the nonzero shifts.  Hence it is enough to
exclude the single support $S$.

\section{A four-shift obstruction over \texorpdfstring{$\F_4$}{F4}}

Reduce $U$ modulo $2$.  Then
\[
 \Z[\omega]/(2)\cong\F_4,
\]
conjugation is $x\mapsto x^2=x^{-1}$ on $\F_4^\times$, every active
coefficient remains nonzero, and $12=0$.  Write
$x_j\in\F_4^\times$ for the value at $j\in S$ and abbreviate
$x_a/x_b=x_a\overline{x_b}$.  The support intersections at shifts
$2,4,8,10$ give
\[
\begin{aligned}
C_2={}&\frac{x_0}{x_2}+\frac{x_1}{x_3}
      +\frac{x_{10}}{x_{12}}+\frac{x_{26}}{x_{28}},\\
C_4={}&\frac{x_3}{x_7}+\frac{x_{12}}{x_{16}}
      +\frac{x_{22}}{x_{26}},\\
C_8={}&\frac{x_2}{x_{10}}+\frac{x_{28}}{x_1},\\
C_{10}={}&\frac{x_0}{x_{10}}+\frac{x_2}{x_{12}}
      +\frac{x_{12}}{x_{22}}+\frac{x_{16}}{x_{26}}
      +\frac{x_{26}}{x_1}+\frac{x_{28}}{x_3}.
\end{aligned}
\]

Suppose all four vanish.  Because the field has characteristic two,
$C_8=0$ says its two nonzero summands are equal.  Put
\[
 p=\frac{x_2}{x_{10}}=\frac{x_{28}}{x_1}.
\]
Multiplication by $p$ gives
\[
 pC_2=
 \frac{x_0}{x_{10}}+\frac{x_{28}}{x_3}
 +\frac{x_2}{x_{12}}+\frac{x_{26}}{x_1},
\]
the sum of four terms of $C_{10}$.  Hence $C_2=C_{10}=0$ forces
\[
 \frac{x_{12}}{x_{22}}=\frac{x_{16}}{x_{26}},
 \qquad\text{and therefore}\qquad
 \frac{x_{12}}{x_{16}}=\frac{x_{22}}{x_{26}}.
\]
The last two summands of $C_4$ are equal and cancel, leaving
\[
 C_4=\frac{x_3}{x_7}\ne0,
\]
a contradiction.

\begin{theorem}\label{thm:cgw}
No circulant $\CGW(35,12;6)$ exists.
\end{theorem}

\begin{proof}
The characteristic-three reduction of a hypothetical matrix has one of the
420 enumerated supports.  Lemma \ref{lem:affine-orbit} moves it to $S$, where
the characteristic-two reduction contradicts the four displayed
correlations.
\end{proof}

\section{The \texorpdfstring{$\CW(105,36)$}{CW(105,36)} corollary}

We import exactly one multiplier result.  Theorem~4.1 of
Arasu--Gordon--Zhang~\cite{ArasuGordonZhang2021}, attributed there to
McFarland, states that if $M\in\operatorname{ICW}_d(m,k)$,
$\gcd(m,k)=1$, and an integer $t$ is congruent modulo $m$ to a power of
each prime divisor of $k$, then $t$ is a multiplier of $M$.  This theorem is
an external dependency, not a computation in the present work.

For a hypothetical $\CW(105,36)$, write
$A(X)=\sum_{j=0}^{104}a_jX^j$.  Its three-fold contraction is defined
coefficientwise by
\[
 b_q=\sum_{r=0}^{2}a_{q+35r},\qquad q\in\mathbb Z_{35}.
\]
Then $B(X)=\sum_{q=0}^{34}b_qX^q$ is an
$\operatorname{ICW}_3(35,36)$, as in the contraction preceding Lemma~2.5
of the cited paper.  Every hypothesis of Theorem~4.1 is explicit:
\[
 \gcd(35,36)=1,\qquad 36=2^2\,3^2,\qquad
 4\equiv2^2\pmod {35},\qquad 4\equiv 3^{10}\pmod {35}.
\]
Thus $4$ is a multiplier of the contraction.  Since
$\gcd(4-1,35)=1$, a translate is fixed by $4$.  The complete finite
enumeration of these fixed integer rows, independently reproduced by two
Python enumerators and a standalone C++ sweep, leaves exactly
\[
\begin{aligned}
 B_1(X)&=-3+3X^5+3X^{10}+3X^{20},\\
 B_2(X)&=-3+3X^{15}+3X^{25}+3X^{30}.
\end{aligned}
\]
We treat both rows directly; their affine equivalence is not used below.

Use the CRT coordinates $i=3q+35r$, with $q\in\Z_{35}$ and
$r\in\Z_3$.  The two rows become, respectively,
\[
\begin{aligned}
 \widetilde B_1(X)&=-3+3X^{15}+3X^{25}+3X^{30},\\
 \widetilde B_2(X)&=-3+3X^5+3X^{10}+3X^{20}.
\end{aligned}
\]
For either row, the unique column of sum $-3$ is $(-1,-1,-1)$ and the
three columns of sum $3$ are $(1,1,1)$.  A column of sum zero is either
zero or one of the six permutations of $(1,-1,0)$.  The four constant
columns contribute 12 to the weight; the remaining weight is 24, and each
active zero-sum column contributes 2.  Hence either contraction row forces
exactly twelve active columns.

For a column $(a_0,a_1,a_2)$ take its nontrivial fibre character
\[
 D_q=a_0+a_1\omega+a_2\omega^2.
\]
The six possible active-column patterns are
\[
\begin{array}{c|cccccc}
e&0&1&2&3&4&5\\ \hline
(a_0,a_1,a_2)&(1,-1,0)&(1,0,-1)&(0,1,-1)&
(-1,1,0)&(-1,0,1)&(0,-1,1)
\end{array}
\]
and satisfy $D_q=(1-\omega)\zeta_6^e$.  Thus
$D=(1-\omega)U$ for a length-35 sequence $U$ with twelve nonzero sixth
roots; constant and zero columns have character value zero.

Here is the projection coefficient by coefficient.  Put
\[
 R(t,s)=\sum_{q\in\Z_{35}}\sum_{r\in\Z_3}
 a_{q,r}a_{q+t,r+s}.
\]
The CRT map is a group isomorphism, so the $\CW(105,36)$ equations say
$R(0,0)=36$ and $R(t,s)=0$ for every other pair.  Substituting
$u=r+s$ gives
\[
\begin{aligned}
 \sum_qD_q\overline{D_{q+t}}
 &=\sum_{q,r,u}a_{q,r}a_{q+t,u}\omega^{r-u}\\
 &=\sum_{s=0}^2\omega^{-s}R(t,s)
 =\begin{cases}36,&t=0,\\0,&t\ne0.\end{cases}
\end{aligned}
\]
Since $D=(1-\omega)U$ and
$(1-\omega)(1-\overline{\omega})=3$, the preceding identity gives
\[
 DD^\#=3UU^\#=36.
\]
The group ring $\mathbb Z[\omega][C_{35}]$ is free as an abelian group,
so multiplication by $3$ is injective and the factor cancels
coefficientwise.  Hence $UU^\#=12$.
This would be a circulant $\CGW(35,12;6)$, contrary to
Theorem~\ref{thm:cgw}.

\begin{theorem}
No circulant weighing matrix $\CW(105,36)$ exists.
\end{theorem}

\section{Reproducibility and evidence boundary}

The proof-bearing computation is the complete enumeration of two ternary
$[35,12]$ codes and the equality of their 420 supports with one affine orbit.
Python and standalone C++ implementations regenerate the complete ordered
certificate independently.  SageMath 10.9 and GAP 4.12.1 implementations
reproduce the classification by distinct CAS routes.  These are independent
implementations within the present project, not independent human
authorship.

A separate C++ checker ignores the affine reduction and checks all
\[
 420\cdot3^{11}=74{,}401{,}740
\]
global-phase-normalized $\F_4$ assignments, finding no solution.  This is
redundant validation, not a logical dependency of the four-shift argument.
An independent standard-library audit checks the factor branches, explicit
word, support certificate, affine stabilizer, four correlation term lists,
fibre table, contraction convention, and the publication manifest.

All theorem decisions use exact integer or finite-field arithmetic; no
incomplete or non-exact computation is used.  The imported contraction
theorem and the apparently new
support-preserving two-characteristic obstruction both remain subject to
specialist review.  Related published work and the maintained matrix
repository are cited in~\cite{GordonRepository2026,Gordon2026}; no priority
claim is made beyond those sources.

\section*{Declaration of generative AI and AI-assisted technologies}

During the development of this work, the author used OpenAI GPT-5.6 Sol,
accessed through ChatGPT and Codex, as a substantive interactive research and
software-development assistant.  The model assisted in exploring candidate
algebraic reductions, formulating and stress-testing proof steps, generating
and reviewing exact verification code, designing adversarial checks, and
organizing and editing the manuscript.  Its outputs were treated as proposals,
not as mathematical evidence or authoritative sources.  The author reviewed
all retained arguments, code, computations, and citations; they were accepted
only after exact checks, cross-implementation comparison, or direct source
verification.  The author takes full
responsibility for the mathematical claims, software, citations, and final
manuscript.

\bibliographystyle{plain}
\bibliography{references}

\end{document}
